% ===================================================================
% Supplemental Material
% Have Optical Wigner's Friend Experiments Been Blind to a Geometric
% Degree of Freedom?
% ===================================================================

\documentclass[12pt]{article}
\usepackage[utf8]{inputenc}
\usepackage{amsmath,amssymb,graphicx,booktabs,hyperref,geometry}
\geometry{margin=1in}

\begin{document}

\title{Supplemental Material:\\
Have Optical Wigner's Friend Experiments Been Blind to a Geometric Degree of Freedom?}
\author{VietVunVut (Viet~--~Nguyen~Xuan)}
\email{viet@vvvqmrf.com}
\homepage{https://vvvqmrf.com}
\affiliation{Independent Researcher, Vietnam}
\date{2026-05-26}
\maketitle

\renewcommand{\thesection}{S\arabic{section}}
\renewcommand{\theequation}{S\arabic{equation}}
\renewcommand{\thetable}{S\arabic{table}}

% ===================================================================
\section{Literature Search and Algebraic Proof}
\label{sec:S1}

\subsection{Full Proof of Equatorial Cancellation Theorem}

\textbf{Theorem.} Let Friend $F$ measure in $z$-basis and Superobserver $W$ at Bloch angles $(\theta,\phi)$.
With $f_\perp(b,d) = 1 - |\langle b|d\rangle|^2$:
\begin{equation}
f_\perp(+1, H) - f_\perp(-1, H) = -\cos\theta
\end{equation}
Vanishes iff $\theta = \pi/2$.

\subsubsection*{Step 1: $W$'s measurement basis}
\begin{align}
|b=+1\rangle &= \cos(\theta/2)|H\rangle + e^{i\phi}\sin(\theta/2)|V\rangle,\\
|b=-1\rangle &= \sin(\theta/2)|H\rangle - e^{i\phi}\cos(\theta/2)|V\rangle.
\end{align}

\subsubsection*{Step 2: Overlaps with $F$'s outcomes}
\begin{align}
|\langle b=+1|H\rangle|^2 &= \cos^2(\theta/2), \quad |\langle b=+1|V\rangle|^2 = \sin^2(\theta/2),\\
|\langle b=-1|H\rangle|^2 &= \sin^2(\theta/2), \quad |\langle b=-1|V\rangle|^2 = \cos^2(\theta/2).
\end{align}
$\phi$ drops out: $|e^{i\phi}|^2 = 1$. Overlaps depend \emph{only} on $\theta$.

\subsubsection*{Step 3: $f_\perp$ values}
\begin{align}
f_\perp(+1,H) &= 1 - \cos^2(\theta/2) = \sin^2(\theta/2),\\
f_\perp(-1,H) &= 1 - \sin^2(\theta/2) = \cos^2(\theta/2),\\
f_\perp(+1,V) &= 1 - \sin^2(\theta/2) = \cos^2(\theta/2),\\
f_\perp(-1,V) &= 1 - \cos^2(\theta/2) = \sin^2(\theta/2).
\end{align}

\subsubsection*{Step 4: Outcome-dependence parameter}
\begin{align}
f_\perp(+1, H) - f_\perp(-1, H) &= \sin^2(\theta/2) - \cos^2(\theta/2)\\
&= -( \cos^2(\theta/2) - \sin^2(\theta/2))\\
&= -\cos\theta.
\end{align}
Vanishes iff $\theta = \pi/2$. At this angle, all four $f_\perp = 1/2$, constant across all outcome pairs. $\square$

\subsubsection*{Generality}
At $\theta = \pi/2$, $|\langle b|d\rangle|^2 = 1/2$ for all $b,d$. Any $g(|\langle b|d\rangle|^2)$ therefore takes the same value for all outcome pairs. The equatorial plane is a fixed point for the entire class --- the cancellation is not specific to $f_\perp$.

\subsection{BSM Equivalence for Proietti et al.\ (2019)}

Proietti et al.\ use Bell-state measurement (BSM) rather than projective measurement. BSM projects onto four Bell states $\{|\Phi^+\rangle, |\Phi^-\rangle, |\Psi^+\rangle, |\Psi^-\rangle\}$. For the singlet-state source, each Bell outcome occurs with equal probability $1/4$, and conditioned on any Bell outcome $b$, the Friend's recorded outcome $d$ is equally likely to be $H$ or $V$. Consequently $|\langle b|d\rangle|^2 = 1/2$ for all $b,d$, identical to the equatorial condition $\theta = \pi/2$ derived in the main text.

\subsection{Literature Search Methodology}

\textbf{Search date:} 2026-05-15 to 2026-05-24.

\textbf{Databases:} Google Scholar, arXiv (quant-ph), Web of Science, InspireHEP.

\textbf{Search strings} (Boolean combinations):
\begin{itemize}
\item \textbf{Group A (subject):} ``Wigner's friend'' OR ``extended Wigner'' OR ``Frauchiger-Renner'' OR ``Local Friendliness''
\item \textbf{Group B (parameter):} ``equatorial measurement'' OR ``Bloch sphere polar angle'' OR ``outcome dependence'' OR ``geometric constraint'' OR ``measurement basis angle''
\end{itemize}

\textbf{Inclusion criteria:}
\begin{enumerate}
\item Optical EWF implementation with Friend$+$Superobserver structure
\item Published in peer-reviewed venue or arXiv
\item Reports measurement settings from which polar angle $\theta$ can be determined
\end{enumerate}

\textbf{Results:} $\sim$200 titles screened $\to$ 47 full-text examined $\to$ 2 published optical EWF experiments identified. Both operate at equatorial geometry ($\theta = \pi/2$ or equivalent via BSM). No published experiment varies $\theta$ from $\pi/2$ for any purpose.

\textbf{Additional sources examined:} Supplemental Material of Bong et al.\ (2020); Methods and Supplementary Information of Proietti et al.\ (2019); LF derivations in Frauchiger-Renner (2018) and Wiseman-Cavalcanti-Rieffel (2023); Bell/LF review by Brunner et al.\ (2014); multipartite, sequential, and possibilistic LF extensions.

\textbf{Limitation:} We cannot rule out unpublished results or implementations outside our database scope. The survey establishes that within the indexed, peer-reviewed literature, $\theta$ has not been varied.

% ===================================================================
\section{Numerical Methods and Statistical Robustness}
\label{sec:S2}

\subsection{Exact Numerical Integration Method}

All numerical predictions use the density matrix for the singlet state with finite visibility:
\begin{equation}
\rho_\mu = \mu|\Phi^-\rangle\langle\Phi^-| + \frac{1-\mu}{4}I.
\end{equation}

For the overlap-dependent model Eq.~(2--3) of the main text:
\begin{equation}
P(a,b \mid x,y) = P_{\rm QM}(a,b \mid x,y) \cdot [1 - \beta f_\perp(b,d)] / Z,
\end{equation}
with full renormalization $Z = \sum_{a,b} P_{\rm QM}(a,b \mid x,y) \cdot [1 - \beta f_\perp(b,d)]$.

All computations performed via \texttt{K9S12\_proposal.py}, \texttt{statistical\_significance.py}. The algebraic cancellation theorem is independently verified via symbolic computation (\texttt{universal\_theorem\_lf\_check.py}).

\subsection{Complete Correlator Table ($\theta = 31^\circ$, $\mu = 0.95$)}

\begin{table}[h]
\centering
\caption{All nine $\langle A_x B_y\rangle$ correlators with uncertainties ($N = 91{,}000$).}
\begin{tabular}{ccc}
\toprule
$(x,y)$ & $\langle AB\rangle_{\rm QM}$ & $\sigma$ \\
\midrule
$(1,1)$ & $-1.0000$ & $0.0000$ \\
$(1,2)$ & $-0.8572$ & $0.0017$ \\
$(1,3)$ & $-0.8572$ & $0.0017$ \\
$(2,1)$ & $-0.8572$ & $0.0017$ \\
$(2,2)$ & $-0.5045$ & $0.0029$ \\
$(2,3)$ & $-0.8933$ & $0.0015$ \\
$(3,1)$ & $-0.8572$ & $0.0017$ \\
$(3,2)$ & $-0.8933$ & $0.0015$ \\
$(3,3)$ & $-0.8829$ & $0.0016$ \\
\bottomrule
\end{tabular}
\end{table}

\subsection{Grid Search Optimization}

A grid search over $(\theta, \phi_2, \phi_3, \beta_{\rm Bob})$ maximizes:
\begin{equation}
\text{FOM}(\theta) = \min\left( n_\sigma^{\rm LF}(\theta), n_\sigma^{\rm signal}(\theta, \beta) \right),
\end{equation}
yielding $\theta = 31^\circ$, $\phi_2 = 112^\circ$, $\phi_3 = 217^\circ$, $\beta_{\rm Bob} = 20^\circ$. FOM remains above $5\sigma$ for $\theta \in [20^\circ, 55^\circ]$ at $\beta = 0.07$, $N = 91{,}000$.

\subsection{Monte Carlo Validation}

10{,}000 runs: Gen LF~1 $\ge 5\sigma$ in 99.97\%. $\beta = 0.07$: ${>}99\%$ power (single); $\beta = 0.05$: $\sim$90\% (combined). Conservative Bayesian analysis (20\% inflation): $\beta_{\rm min} \approx 0.046$ (combined).

\subsection{Angular Stability}

LF significance varies ${<}3\%$ across $\Delta\theta = \pm 5^\circ$. $\delta\langle AB\rangle \propto \cos\theta$: $\sim$1\% at $\pm 1^\circ$ (Bong's reported precision).

\subsection{$\phi$-Scramble Control Protocol}

Rotate HWP to $N_\phi \geq 10$ random azimuthal angles $\phi \in [0^\circ, 360^\circ)$, fit $\delta\langle AB\rangle(\phi) = A + B\cos(2\phi) + C\sin(2\phi)$. Test: $A \neq 0$ with $B, C \approx 0$ indicates geometric $\cos\theta$ signal; $B$ or $C \neq 0$ indicates birefringence/alignment artifacts.

% ===================================================================
\section{Interpretations and Physical Context}
\label{sec:S3}

\subsection{Relation to Contextuality}

Eq.~(2) of the main text is not a contextual hidden-variable model. There is no unobserved $\lambda$; the Friend outcome $d$ is an observed macroscopic record. Dependence is on registration geometry ($|\langle b|d\rangle|^2$) rather than on which observables are co-measured.

\begin{table}[h]
\centering
\caption{Comparison of dependence frameworks.}
\begin{tabular}{p{0.25\textwidth}p{0.18\textwidth}p{0.25\textwidth}p{0.25\textwidth}}
\toprule
Framework & Hidden variable? & Dependence on & Status of $d$ \\
\midrule
KS Contextuality & Yes ($\lambda$) & Measurement context & Unobserved \\
Standard contextuality & No & Co-measured observables & N/A \\
Retrocausal models & Varies & Future choices & Varies \\
Eq.~(2) registration & No & $|\langle b|d\rangle|^2$ geometry & Observed record \\
\bottomrule
\end{tabular}
\end{table}

\subsection{Physical Picture}

When a Friend measures in the $z$-basis, they produce a macroscopic record with a definite orientation on the Bloch sphere. A Superobserver measuring at $(\theta,\phi)$ probes this record at a relative angle $|\langle b|d\rangle|^2$. Standard QM assumes that once the Friend's outcome is registered, subsequent measurements factorize. Eq.~(2) parametrizes a possible residual dependence on this geometric relationship.

\subsection{Quantum Interpretations}

\textbf{Copenhagen.} Silent on registration-layer structure --- the experiment tests a structural degree of freedom that no standard interpretation addresses.

\textbf{Many-Worlds (Everett).} $\beta \neq 0$ would provide the first quantitative probe of inter-branch structural correlations; a null result constrains such effects to $\beta < 0.04$.

\textbf{Relational QM (Rovelli).} Tests whether relational outcomes leave geometry-dependent measurable traces --- a signature current RQM formulations neither predict nor preclude.

\textbf{Objective Collapse.} CSL and Diosi-Penrose models modify dynamics, not registration geometry. The overlap-dependent class is structurally orthogonal.

\textbf{QBism.} Parametrizes whether agent experiences exhibit structural coupling across agents at the registration layer. A null result constrains agent-agent coupling at $\sim\!10^{-2}$.

\subsection{General Probabilistic Theories (GPT) Context}

Within the GPT framework, the overlap-only deformation defines a one-parameter family of neighboring GPTs deviating from the quantum point along the registration-geometry direction. This is a structural degree of freedom orthogonal to standard nonlocality.

\subsection{Weak Measurement and Postselection}

In weak measurement, the weak value $A_w = \langle \psi_f | A | \psi_i \rangle / \langle \psi_f | \psi_i \rangle$ depends explicitly on the pre-postselection overlap. The basis-alignment dependence in Eq.~(2) shares the structural feature of geometric coupling between measurement records at different observational levels. The analogy is structural, not dynamical.

\subsection{Scale Comparison}

\begin{table}[h]
\centering
\caption{Phenomenological parameter scales.}
\begin{tabular}{lcc}
\toprule
Parameter class & Typical constraint & Reference \\
\midrule
Photon-sector SME & ${<}10^{-23}$ & Colladay \& Kosteleck\'y (1997) \\
CSL collapse rate $\lambda$ & ${<}10^{-16}$ s${}^{-1}$ & Various \\
Weak-measurement anomalies & $\sim\!10^{-2}$ & Aharonov et al.\ (1988) \\
Overlap-dependent $\beta$ (this work) & $\sim\!10^{-2}$ (proposed) & --- \\
\bottomrule
\end{tabular}
\end{table}

\subsection{Multi-Observer Extension (Preliminary)}

For $N > 2$ observers, the geometric cancellation structure generalizes. Preliminary analysis suggests $\sim\!11\times$ amplification at $\beta = 0.3$ for three observers, conditional on bridge theorems connecting registration-layer structure to QM observables for $N > 2$. Rigorous derivation is left for future work.

\end{document}
